\documentclass[11pt]{article}

\usepackage[margin=1in]{geometry}
\usepackage{amsmath}
\usepackage{amssymb}
\usepackage{booktabs}
\usepackage{siunitx}
\usepackage{hyperref}
\usepackage{xcolor}

\hypersetup{
  colorlinks=true,
  linkcolor=blue!50!black,
  urlcolor=blue!50!black,
}

\title{The Physics Behind the 2\,inch (50\,mm) Main Drain Flow Test}
\author{Implementation notes for AXA XL PRC.14.1.2.2}
\date{}

\begin{document}
\maketitle

\section*{What we are doing}

A 2\,inch main drain test estimates the available municipal (or yard) water
supply at a sprinkler riser, without flowing a hydrant. We open the
2\,inch drain wide, let the system stabilize, and read \emph{two} pressures
on the riser gauge:
\begin{itemize}
  \item the \textbf{static pressure} $P_S$ before opening the drain, and
  \item the \textbf{residual pressure} $P_R$ while the drain is flowing.
\end{itemize}

What we want is the \textbf{drain flow} $Q$ in gallons per minute (gpm).
We can't measure $Q$ directly without a pitot, but we can compute it from
$P_R$ if we know the equivalent length of the drain piping. The chart in
PRC.14.1.2.2 Figure~1 is a graphical solution of the same equations
derived below.

\section{The drain as a hydraulic circuit}

While the drain is flowing, the only thing between the riser gauge and the
atmosphere is a length of 2\,inch pipe with some fittings. Energy is lost
in exactly two places:

\begin{enumerate}
  \item \textbf{Outlet pressure} ($P_{\text{out}}$): the kinetic energy
        the water carries as it leaves the discharge. This is the velocity
        head, captured by the discharge coefficient of the outlet.
  \item \textbf{Friction loss} ($P_{\text{f}}$): pressure lost rubbing
        against the pipe walls and through fittings between the riser and
        the outlet.
\end{enumerate}

The elevation difference between the riser gauge and the outlet is usually
a couple of feet, contributing $<1$\,psi, and is ignored. Velocity
pressure at the gauge is also ignored on the assumption that the gauge is
tapped into the riser, not into the drain line (this is one of the
stated parameters of the method).

Conservation of energy then gives the working equation:
\begin{equation}
\boxed{\;P_R \;=\; P_{\text{out}}(Q) \;+\; P_{\text{f}}(Q, L)\;}
\label{eq:main}
\end{equation}
where $L$ is the equivalent length of 2\,inch pipe between the riser and
the outlet.

\section{Outlet pressure term}

The standard fire-protection discharge equation for a smooth orifice or
short tube is
\begin{equation}
Q \;=\; 29.84\, c\, d^{2}\, \sqrt{P_{\text{out}}}
\label{eq:disch}
\end{equation}
with $Q$ in gpm, $d$ the orifice diameter in inches, $P_{\text{out}}$ in psi,
and $c$ a dimensionless discharge coefficient.

PRC.14.1.2.2 reports measured outlet coefficients of approximately
$c=0.90$ for a reasonably well-reamed 2\,inch outlet and $c=0.85$ for an
unreamed 2\,inch outlet or for discharge through a 45$^\circ$ elbow. The
guideline adopts the more conservative
\[
c \;=\; 0.85
\]
and we follow suit. With $d = 2$\,inch:
\begin{equation}
Q \;=\; 29.84\,(0.85)\,(2)^{2}\,\sqrt{P_{\text{out}}}
   \;=\; 101.456\,\sqrt{P_{\text{out}}},
\end{equation}
which we invert for use in Eq.~\eqref{eq:main}:
\begin{equation}
\boxed{\;P_{\text{out}}(Q) \;=\; \left(\dfrac{Q}{101.456}\right)^{2}\;}
\label{eq:pout}
\end{equation}

\section{Friction loss term (Hazen--Williams)}

Friction loss in the 2\,inch pipe is computed with the Hazen--Williams
equation, the standard tool in fire-protection hydraulics. In US units,
pressure loss per unit length is
\begin{equation}
p_{f} \;=\; \dfrac{4.52\, Q^{1.852}}{C^{1.852}\, d^{4.87}}
\qquad [\text{psi/ft}]
\label{eq:hw}
\end{equation}
with $Q$ in gpm, $d$ the \emph{actual inside diameter} in inches, and $C$
the Hazen--Williams roughness coefficient. PRC.14.1.2.2 specifies $C=120$
(typical for steel sprinkler pipe).

The pipe is nominal 2\,inch, Schedule~40. The actual inside diameter is
$d = 2.067$\,inch, not 2.000\,inch. (Using the nominal value would
under-predict friction loss by $\sim 18\%$ at fixed flow.)

Total friction loss over the equivalent length $L$ is
\begin{equation}
\boxed{\;P_{\text{f}}(Q, L) \;=\; L \cdot \dfrac{4.52\, Q^{1.852}}{120^{1.852}\,(2.067)^{4.87}}\;}
\label{eq:pf}
\end{equation}

Plugging in the constants:
\[
120^{1.852} \;\approx\; 7\,096, \qquad (2.067)^{4.87} \;\approx\; 36.96,
\]
\[
\dfrac{4.52}{7\,096 \cdot 36.96} \;\approx\; 1.724 \times 10^{-5}
\quad [\text{psi/ft per } Q^{1.852}].
\]
So in compact form,
\begin{equation}
P_{\text{f}}(Q, L) \;\approx\; 1.724 \times 10^{-5}\, L\, Q^{1.852}
\qquad [\text{psi}].
\label{eq:pf-numeric}
\end{equation}

\section{Equivalent length $L$}

$L$ is the length of 2\,inch \emph{straight pipe} that would produce the
same friction loss as the actual run, including fittings. Each fitting is
assigned an equivalent length per Table~1 of PRC.14.1.2.2:

\begin{center}
\begin{tabular}{lc}
\toprule
Fitting & Equivalent length (ft) \\
\midrule
Angle valve            & 29 \\
Globe valve            & 58 \\
Gate valve             & 1  \\
$90^\circ$ elbow       & 5  \\
$45^\circ$ elbow       & 2  \\
Tee (flow turns)       & 10 \\
Cross (flow turns)     & 10 \\
\bottomrule
\end{tabular}
\end{center}

So if the drain run is $\ell$ feet of straight pipe plus $n_i$ of each
fitting type with equivalent length $L_i$,
\begin{equation}
L \;=\; \ell \;+\; \sum_{i}\, n_i\, L_i.
\end{equation}

\section{Combined equation}

Substituting \eqref{eq:pout} and \eqref{eq:pf} into \eqref{eq:main}:
\begin{equation}
\boxed{\;
P_R(Q, L) \;=\; \left(\dfrac{Q}{101.456}\right)^{2}
\;+\; \dfrac{4.52\, L\, Q^{1.852}}{120^{1.852}\,(2.067)^{4.87}}
\;}
\label{eq:full}
\end{equation}

This is one equation in one unknown ($Q$). The chart in Figure~1 of
PRC.14.1.2.2 is a graphical plot of Eq.~\eqref{eq:full} for several values
of $L$.

\section{The inverse problem (going from $P_R$ to $Q$)}

In practice we measure $P_R$ and want $Q$. Eq.~\eqref{eq:full} is not
algebraically solvable for $Q$ because $Q$ appears with two different
exponents ($Q^2$ and $Q^{1.852}$). Two ways to handle this:

\paragraph{Graphical (the guideline's method).} Read horizontally across
Figure~1 at $P_R$, find the curve labelled with your $L$ (interpolate
between curves if needed), then read down to $Q$. Easy and visual;
limited to about $\pm5\%$ by the resolution of the chart.

\paragraph{Numerical (what the web tool does).} Since
$P_R(Q, L)$ is strictly increasing in $Q$ for $Q>0$, bisection converges
unconditionally. Pick $Q_{\text{lo}}=0$ and $Q_{\text{hi}}$ generous
(2000~gpm is plenty). At each iteration evaluate the midpoint; if the
computed $P_R$ is less than the measured value, the true $Q$ must be
higher (replace $Q_{\text{lo}}$ with the midpoint); otherwise replace
$Q_{\text{hi}}$. After $n$ iterations the interval is $2000/2^{n}$ gpm
wide, so $\sim 80$ iterations gives essentially exact answers.

\section{Worked example (from PRC.14.1.2.2)}

A plant has two risers, A and B, with the drain piping described below.
A static of $P_S = 100$\,psi is recorded at riser A. Two flow scenarios
are run.

\subsection*{Equivalent lengths}

\begin{center}
\begin{tabular}{lcc}
\toprule
Component & Drain A & Drain B \\
\midrule
2\,inch pipe (ft)        & 8                & 22 \\
Angle valves $\times 29$ & $1 \cdot 29 = 29$ & $1 \cdot 29 = 29$ \\
$90^\circ$ ells $\times 5$ & $1 \cdot 5 = 5$   & $3 \cdot 5 = 15$ \\
$45^\circ$ ells $\times 2$ & $1 \cdot 2 = 2$   & $1 \cdot 2 = 2$  \\
\midrule
\textbf{Total $L$ (ft)}  & \textbf{44}      & \textbf{68} \\
\bottomrule
\end{tabular}
\end{center}

\subsection*{Scenario 1 --- Drain A alone}

Measured residual at riser A: $P_R = 86$\,psi. We solve
Eq.~\eqref{eq:full} for $Q$ with $L = 44$\,ft.

Quick consistency check at $Q = 450$\,gpm:
\begin{align*}
P_{\text{out}} &= \left(\dfrac{450}{101.456}\right)^{2}
                = (4.435)^{2}
                = 19.7~\text{psi} \\[4pt]
Q^{1.852} &= 450^{1.852} \approx 82{,}000 \\[4pt]
P_{\text{f}}  &= 1.724\times10^{-5} \cdot 44 \cdot 82{,}000
                \approx 62.2~\text{psi}
                \quad\text{(more precise: } 66.3) \\[4pt]
P_R           &= 19.7 + 66.3 \;\approx\; 86~\text{psi.} \quad\checkmark
\end{align*}
The bisection solver lands on $Q \approx 448$~gpm; PRC.14.1.2.2 reports
$450$~gpm.

\subsection*{Scenario 2 --- Drains A and B flowing simultaneously}

Now both drains are open. The riser gauges read $P_R^{A} = 64$\,psi at
riser A and $P_R^{B} = 70$\,psi at riser B. We solve
Eq.~\eqref{eq:full} \emph{independently for each drain}:

\begin{align*}
\text{Drain A:}\quad & L = 44, \; P_R = 64 \;\Rightarrow\; Q_A \approx 383~\text{gpm} \quad (\text{doc: } 390) \\[2pt]
\text{Drain B:}\quad & L = 68, \; P_R = 70 \;\Rightarrow\; Q_B \approx 333~\text{gpm} \quad (\text{doc: } 335)
\end{align*}

Total flow during the simultaneous test:
\[
Q_{\text{tot}} \;=\; Q_A + Q_B \;\approx\; 716~\text{gpm}
\qquad (\text{doc: } 725~\text{gpm}).
\]

\subsection*{Two points on the water supply curve}

Using \emph{riser A} as the reference gauge for both scenarios, we have:
\begin{center}
\begin{tabular}{lccc}
\toprule
 & Static $P_S$ & Residual $P_R^A$ & Total flow $Q_{\text{tot}}$ \\
\midrule
Scenario 1 & 100 psi & 86 psi & 450 gpm \\
Scenario 2 & 100 psi & 64 psi & 725 gpm \\
\bottomrule
\end{tabular}
\end{center}

These two points plus the static $(0, P_S)$ define a straight line on
$N^{1.85}$ semi-exponential graph paper, which is the standard
fire-protection format for plotting a water-supply curve.

\section{The water supply curve}

A municipal supply curve on $N^{1.85}$ paper takes the form
\begin{equation}
P(Q) \;=\; P_S \;-\; k\, Q^{1.85},
\label{eq:supply}
\end{equation}
which is linear in $Q^{1.85}$. With two or more $(Q_i, P_i)$ test points
we estimate $k$ by averaging
\[
k_i \;=\; \dfrac{P_S - P_i}{Q_i^{1.85}}.
\]

A common reportable number is the flow at 20\,psi residual:
\begin{equation}
Q_{20} \;=\; \left(\dfrac{P_S - 20}{k}\right)^{1/1.85}.
\end{equation}
This is what would be available if the supply were stretched until the
residual at the test point dropped to the minimum useful sprinkler
pressure of 20\,psi --- a standard way to characterize the supply.

\section{Why this method is conservative}

PRC.14.1.2.2 emphasizes that a drain test result is \emph{typically more
conservative} than a true hydrant flow test. Two reasons:

\begin{enumerate}
  \item At the flows developed in a drain test (a few hundred gpm), an
        alarm check valve or other system check valve may only
        partially open and ``flutter.'' The friction loss across the
        check valve is then significant and unknown, and gets lumped
        into the residual reading. At higher hydrant-test flows the
        valve is fully open and that loss is small.
  \item We adopt $c=0.85$ rather than $0.90$ for the discharge
        coefficient, which under-counts the flow at a given residual.
\end{enumerate}

Both effects bias the inferred supply downwards. That's why the
guideline restricts this method to \textbf{Light Hazard} and
\textbf{Ordinary Hazard} occupancies and refuses to let you extrapolate
the result to other risers or larger flows: the bias is at the test
point only, and assuming it transfers is unsafe.

\section{Numerical constants summary}

\begin{center}
\begin{tabular}{lll}
\toprule
Symbol & Value & Description \\
\midrule
$c$         & $0.85$            & outlet discharge coefficient (conservative) \\
$d_{\text{out}}$ & $2.000$ in   & nominal outlet diameter \\
$d_{\text{pipe}}$ & $2.067$ in  & actual ID of Sch.~40 2\,inch pipe \\
$C$         & $120$             & Hazen--Williams roughness for steel \\
$29.84$     & ---               & US-unit constant in Eq.~\eqref{eq:disch} \\
$4.52$      & ---               & US-unit constant in Eq.~\eqref{eq:hw} \\
$101.456$   & $29.84\,c\,d_{\text{out}}^{2}$ & outlet pre-factor \\
$1.724\times10^{-5}$ & $\dfrac{4.52}{C^{1.852}\,d_{\text{pipe}}^{4.87}}$ & friction pre-factor (per ft, per $Q^{1.852}$) \\
\bottomrule
\end{tabular}
\end{center}

\section*{References}

\begin{itemize}
  \item AXA XL Risk Consulting, \emph{Property Risk Consulting Guidelines
        PRC.14.1.2.2: Water Measurement Using Two Inch (50 mm) Drain
        Tests}, 2020.
\end{itemize}

\end{document}
